\documentclass[11pt,a4paper]{article}
\usepackage[margin=2.5cm]{geometry}
\usepackage[T1]{fontenc}
\usepackage{amsmath}    % align, matrices, \text, \operatorname
\usepackage{amssymb}    % extra symbols, \mathbb, \nexists

% ---------------------------------------------------------------
% MODULE III - Typesetting Mathematics
% Syllabus units 8-11 (Text 1, Tutorial VIII.1 to VIII.8).
% Compiles with a single run of pdflatex.
% ---------------------------------------------------------------

\title{Module III: Typesetting Mathematics}
\author{Study Notes -- MAT5FS112}
\date{}

\begin{document}
\maketitle

% ===============================================================
\section{Basic mathematical equations}
% ===============================================================

\subsection{Inline and displayed}

Put maths between single \verb|$| signs to keep it in the sentence:
the area is $A = \pi r^2$. Use \verb|\[ ... \]| (or the
\texttt{equation*} environment) for a centred display:
\[
  \int_0^1 x^2\,dx = \frac{1}{3}.
\]

\subsection{Numbered equations and labels}

The \texttt{equation} environment numbers the formula. A \verb|\label|
lets you refer to it with \verb|\eqref|:
\begin{equation}
  E = mc^2 \label{eq:einstein}
\end{equation}
Equation~\eqref{eq:einstein} is the mass--energy relation.

\subsection{Scripts, fractions and roots}

\begin{verbatim}
x^{2}          x_{i}          x_{i}^{2}       % superscript, subscript
\frac{a+b}{c}  \sqrt{x}       \sqrt[3]{x}     % fraction, roots
\binom{n}{k}                                  % binomial coefficient
\end{verbatim}

Rendered: $x^{2}$, $x_{i}$, $x_{i}^{2}$, $\frac{a+b}{c}$,
$\sqrt{x}$, $\sqrt[3]{x}$, $\binom{n}{k}$.

Braces group several symbols. Without them, only the next character is
affected: \verb|x^12| gives $x^1 2$, while \verb|x^{12}| gives $x^{12}$.

% ===============================================================
\section{Groups of equations and numbering}
% ===============================================================

\texttt{align} lines up several equations at the \verb|&| character.
Each row ends with \verb|\\|. \texttt{align} numbers every row; the
starred version \texttt{align*} numbers none.

\begin{align}
  (a+b)^2 &= a^2 + 2ab + b^2 \label{eq:sq} \\
  (a-b)^2 &= a^2 - 2ab + b^2 \nonumber
\end{align}

Row~\eqref{eq:sq} is numbered and the second row is not, because
\verb|\nonumber| suppresses its number. A single row with a custom label
uses \verb|\tag{name}|:
\begin{equation}
  \sin^2\theta + \cos^2\theta = 1 \tag{$\star$}
\end{equation}

% ===============================================================
\section{Matrices, dots, delimiters and affixed symbols}
% ===============================================================

\subsection{Matrices}

Inside a display, the \texttt{pmatrix} (round brackets), \texttt{bmatrix}
(square), \texttt{vmatrix} (determinant bars) and \texttt{Bmatrix}
(curly braces) environments set matrices. Rows use \verb|\\| and
columns use \verb|&|.

\[
  A = \begin{pmatrix} 1 & 2 & 3 \\ 0 & -1 & 4 \\ 5 & 6 & 0 \end{pmatrix},
  \qquad
  \det\begin{vmatrix} a & b \\ c & d \end{vmatrix} = ad - bc.
\]

\subsection{Dots}

\[
  x_1, x_2, \ldots, x_n \qquad
  a_{11} + \cdots + a_{nn} \qquad
  \vdots \quad \ddots
\]

Use \verb|\ldots| on the baseline, \verb|\cdots| for centred dots
between operators, \verb|\vdots| and \verb|\ddots| for vertical and
diagonal dots.

\subsection{Delimiters}

Brackets that grow with their content use \verb|\left| and \verb|\right|:
\[
  \left( \frac{x}{y} \right)^{2} \qquad
  \left[ \sum_{k=1}^{n} k \right] \qquad
  \left\{ x \in \mathbb{R} : x > 0 \right\}
\]

Each \verb|\left| needs a matching \verb|\right|; use \verb|\right.|
when you want an invisible closing delimiter.

\subsection{Affixed symbols (accents and decorations)}

\[
  \hat{x} \quad \bar{x} \quad \vec{v} \quad \dot{x} \quad \ddot{x}
  \quad \tilde{a} \quad \overline{AB} \quad \underline{ab}
  \quad \overbrace{a+b+c}^{\text{sum}} \quad \underbrace{x+y}_{\text{pair}}
\]

% ===============================================================
\section{Operators, relations and symbols}
% ===============================================================

\subsection{Greek letters}

Lowercase letters are the name with a backslash: \verb|\alpha|,
\verb|\beta|, \verb|\gamma|, \verb|\theta|, \verb|\lambda|,
\verb|\pi|, \verb|\sigma|, \verb|\omega|. Capital letters are
\verb|\Gamma|, \verb|\Delta|, \verb|\Theta|, \verb|\Lambda|,
\verb|\Sigma|, \verb|\Omega|. Rendered: $\alpha, \beta, \gamma, \theta,
\lambda, \pi, \sigma, \omega$ and $\Gamma, \Delta, \Theta, \Omega$.

\subsection{Big operators and limits}

\[
  \sum_{n=1}^{\infty} \frac{1}{n^2} = \frac{\pi^2}{6}, \qquad
  \prod_{i=1}^{n} i = n!, \qquad
  \lim_{x \to 0} \frac{\sin x}{x} = 1
\]

\subsection{Named functions}

Use \verb|\sin|, \verb|\cos|, \verb|\log|, \verb|\det| and \verb|\max|
rather than italic letters. For a name that LaTeX does not provide,
use \verb|\operatorname{name}| (from \texttt{amsmath}):
$\operatorname{sech}\,x = \dfrac{1}{\cosh x}$.

\subsection{Relations, arrows and set notation}

\[
  x \le y, \quad x \ne y, \quad x \approx y, \quad
  f\colon A \to B, \quad x \mapsto x^2
\]
\[
  a \in A, \quad A \subseteq B, \quad A \cup B, \quad A \cap B, \quad
  \forall\, \varepsilon > 0\ \exists\, \delta > 0
\]

\subsection{Text and font inside maths}

\begin{verbatim}
$\text{if } x > 0$        % text inside maths, uses the document font
$\mathrm{d}x$             % upright d for differentials
$\mathbf{v}$              % bold vector
$\mathbb{R}, \mathbb{N}$  % blackboard bold (amssymb)
$\mathcal{L}$             % calligraphic
\end{verbatim}

Rendered: $\text{if } x > 0$, $\int f\,\mathrm{d}x$, $\mathbf{v}$,
$\mathbb{R}, \mathbb{N}$, $\mathcal{L}$.

% ===============================================================
\section{A worked example combining the tools}
% ===============================================================

Let $A$ be a $2 \times 2$ matrix with eigenvalues $\lambda_1, \lambda_2$.
Then
\begin{align*}
  \operatorname{tr}(A) &= \sum_{i=1}^{2} \lambda_i, \\
  \det(A)              &= \prod_{i=1}^{2} \lambda_i .
\end{align*}

\end{document}
